\section{Overview}

This document examines the sensitivity of the main report's analytical results
to alternative methodological choices. It begins with a data-quality assessment
and then addresses the following analytical areas:

\begin{enumerate}
  \item \textbf{Data quality and missing values} --- a per-series inventory of
    missing observations, their causes, and the treatment applied throughout
    the analysis.
  \item \textbf{Variable selection and coefficient diagnostics} --- full OLS
    regression output for all candidate fair-value models, including coefficient
    estimates, standard errors, \(t\)-statistics, and information criteria
    (AIC/BIC) to support the model selection rationale.
  \item \textbf{Lag structure and cross-correlation analysis} --- autocorrelation
    and partial autocorrelation of CDX IG daily changes, together with
    cross-correlation functions against each candidate driver, to assess whether
    lagged regressors would improve model specification.
  \item \textbf{Regime-switching static fair value (RV vs VIX)} --- a
    sensitivity in which OLS coefficients depend on whether trailing realised
    volatility of the S\&P index (from daily level returns) exceeds VIX by more
    than a fixed margin (percentage points on the spread), with a
    figure showing regime timing.
  \item \textbf{Rolling-window fair value model} --- whether allowing model
    coefficients to vary over time (rolling-window estimation) improves or
    degrades the fair-value specification relative to the static
    global-coefficient approach used in the main report.
\end{enumerate}
